\chapter{Vistas y Proyecciones Ortogonales}

\section{Sistemas de Proyección y Vistas Ortogonales}

La representación de objetos tridimensionales (3D) en un plano
bidimensional (2D) es el desafío central del dibujo técnico. La técnica
para lograr esto con precisión se denomina \textbf{proyección
	ortogonal}. Según la normativa IRAM, se utiliza predominantemente el
sistema europeo (proyección en el Primer Diedro).

\subsection{El concepto de Proyección}

Imaginemos el objeto situado dentro de un cubo de cristal. Las
proyecciones se obtienen proyectando los puntos del objeto sobre las
caras internas de este cubo, siguiendo líneas perpendiculares al plano
de proyección. Al desplegar las caras del cubo sobre un plano, obtenemos
las vistas necesarias para definir el objeto.

\subsection{Vistas Principales}

Las vistas normalizadas permiten describir completamente una pieza:
\begin{itemize}
	\item \textbf{Vista de frente (Alzada):} Es la vista principal que
	      debe mostrar la forma general y las características
	      más importantes de la pieza.
	\item \textbf{Vista superior (Planta):} Se sitúa debajo de la
	      vista de frente.
	\item \textbf{Vista lateral derecha (Perfil):} Se sitúa a la
	      derecha de la vista de frente.
\end{itemize}

Otras vistas incluyen la lateral izquierda, inferior y posterior, aunque
se utilizan solo cuando las tres principales no son suficientes.

\subsection{Correspondencia entre Vistas}

La interdependencia entre estas vistas es absoluta:
\begin{enumerate}
	\item La vista de frente y la vista superior comparten la misma
	      anchura.
	\item La vista de frente y la vista lateral comparten la misma
	      altura.
	\item La vista superior y la vista lateral comparten la misma
	      profundidad.
\end{enumerate}

Esta correspondencia permite verificar la exactitud del dibujo. Las
\textbf{líneas ocultas} (trazos discontinuos) son cruciales para
representar bordes o aristas que no son visibles desde la dirección
de observación, permitiendo comprender la estructura interna de la
pieza sin necesidad de cortes.

\section{Cuestionario}
\begin{enumerate}
	\item ¿En qué se basa el sistema de proyección del Primer Diedro?
	\item Define las vistas principales y su disposición.
	\item ¿Qué es la línea de referencia o línea de tierra?
	\item ¿Por qué son necesarias las líneas discontinuas en una vista?
	\item ¿Qué ocurre si dos vistas no coinciden en sus medidas?
	\item ¿Cómo ayuda el software CAD a proyectar automáticamente estas
	      vistas?
	\item ¿Qué es un sólido en el contexto de CAD?
	\item ¿Por qué es importante elegir correctamente la vista de frente
	      al iniciar un diseño?
\end{enumerate}

\section{Parte Práctica}
\begin{enumerate}
	\item Dibujar las tres vistas principales de un cubo de 50 mm.
	\item Dibujar las vistas de un prisma rectangular de 60x40x20 mm.
	\item Proyectar las vistas de una pieza con forma de L.
	\item Dibujar las vistas de un cilindro apoyado en la base.
	\item Representar una pieza con un agujero central mediante líneas discontinuas.
	\item Dibujar las vistas de un cono truncado.
	\item Proyectar una pieza con caras inclinadas.
	\item Dibujar las vistas de una pieza mecánica simple.
	\item Crear una vista de detalle de un elemento pequeño.
	\item Comprobar la correspondencia entre las tres vistas.
\end{enumerate}
